% Préambule APMEP - Annales
% Auteurs
% tapuscrit : François Kriegk
% relecture :
% figures TikZ : François Kriegk
% figures pstricks : Denis Vergès
% modifications :
% Remerciements :
% version 1 : 2026-07-05
%
\documentclass[12pt,a4paper,french]{article}

% Packages
%%% codage et police utilisée
\usepackage[T1]{fontenc}
\RequirePackage{iftex}
\iftutex % moteur utf-8 lualatex, xelatex
	\usepackage{fourier-otf} % remplace fourier pour moteur utf8 (LuaLaTeX, XeLaTeX)
	% Alternatives nécéssite (LuaLaTeX, XeLaTeX)
    \usepackage{hyperref}
	%\usepackage{fontspec}% pour utiliser des fontes système
\else % pour les autres moteurs latex, pdflatex
	\usepackage[utf8]{inputenc}
	\usepackage{fourier}
	\usepackage{amsmath,amssymb}
    \usepackage[dvips]{hyperref}
    \renewcommand{\ttdefault}{lmtt}
\fi
\usepackage[normalem]{ulem}
\usepackage{pifont}
\usepackage{lscape}
% tableaux
\usepackage{multicol}
\usepackage{multirow}
\usepackage{tabularx}
\usepackage{tabularray}
% texte et composition
\usepackage{textcomp}
\usepackage{babel}
\usepackage[np]{numprint}
%
\usepackage{fancyhdr}
\usepackage{makeidx}
\makeindex
% pour le symbole €
\usepackage{eurosym}
\DeclareRobustCommand{\officialeuro}{%
  \ifmmode\expandafter\text\fi
  {\fontencoding{U}\fontfamily{eurosym}\selectfont e}}
\AtBeginDocument{\renewcommand{\euro}{\officialeuro}}
% couleurs étendues
\usepackage{xcolor}
% pour les listes
\usepackage{enumitem}
% Enumérations réglage par défaut
\setlist[itemize]{label=\textbullet,leftmargin=5mm,labelwidth=3mm,labelsep=2mm,align=left}
\setlist[enumerate]{labelindent=0pt, leftmargin=5mm, labelsep=2mm, labelwidth=3mm,labelsep=2mm, align= left}
\setlist[enumerate,1]{label = \textbf{\arabic*.}}
\setlist[enumerate,2]{label = \textbf{\alph*.}}
\setlist[enumerate,3]{label = \textbf{\roman*.}}
% pour l'ajustement des boites (figures…)
\usepackage{adjustbox}
% pour ne pas commencer un exercice en fin de page
\usepackage{needspace}
%
% Taille de la zone d'écriture
\usepackage[left=1.5cm,right=1.5cm,top=3cm,bottom=3cm]{geometry}
\setlength{\headheight}{15pt}
\setlength{\parindent}{0mm}
%
% Formatage du pdf
\hypersetup{%
    pdfauthor = {APMEP},
    pdfsubject = {Troisième, Série générale},
    pdftitle = {Brevet des collèges    -- Session 2026},
    allbordercolors = white,
    pdfstartview=FitH
}
% Pour scratch (code)
\usepackage{scratch3}
% Pour le code Python
\usepackage{listings}
\lstdefinestyle{python}{
  language=Python,
  basicstyle=\ttfamily\small,
  keywordstyle=\color{blue}\bfseries,
  commentstyle=\color{gray}\itshape,
  stringstyle=\color{orange},
  numbers=left,
  numberstyle=\tiny\color{gray},
  breaklines=true,
  frame=single,
  backgroundcolor=\color{gray!5},
  tabsize=4
}
  % usage (exemple)
  %\begin{lstlisting}[style=python]
%def seuil():
%    n = 1
%    p = 0.5
%    while p < 0.6 :
%       n = n + 1
%       p = p*7/15 + 1/3
%    return n
%\end{lstlisting}


% Pour une visualisation de type tableur
\usepackage{pas-tableur}
%
%% pour PsTricks
\usepackage{pst-bezier}%% AVANT les autres PST
\usepackage{pstricks-add}
\usepackage{pst-func,pst-tree,pst-circ}
\usepackage{pst-eucl}
%% pour TikZ
\usepackage{tikz,pgfplots,pgf,tkz-tab,tikzpagenodes}
\usetikzlibrary{calc,patterns,decorations,arrows.meta,patterns.meta,arrows}
\pgfplotsset{compat=1.18,/pgf/number format/.cd,use comma}
\usepackage{icomma} % gestion des espaces autour des virgules
\DeclareMathSymbol{;}{\mathbin}{operators}{"3B}% pour un espacement correct du ; en mode maths
%%%%%%%%%%%%%%%%%%%%%%%%%%%%%%%%%%%%%%%%%%
% Définitions
%
% notation en mode mathématique et hors mode mathématique
\newcommand{\R}{\ensuremath{\mathbb{R}}}
\newcommand{\N}{\ensuremath{\mathbb{N}}}
\newcommand{\D}{\ensuremath{\mathbb{D}}}
\newcommand{\Z}{\ensuremath{\mathbb{Z}}}
\newcommand{\Q}{\ensuremath{\mathbb{Q}}}
\newcommand{\C}{\ensuremath{\mathbb{C}}}
%
\providecommand{\up}{\textsuperscript}
%
\newcommand{\vect}[1]{\overrightarrow{\,\mathstrut#1\,}}
\newcommand{\vectt}[1]{\overrightarrow{\,\mathstrut\text{#1}\,}}
%
\def\Oij{$\left(\text{O};\vect{\imath},\vect{\jmath}\right)$}
\def\Oijk{$\left(\text{O};\vect{\imath},\vect{\jmath},\vect{k}\right)$}
\def\Ouv{$\left(\text{O};\vect{u},\vect{v}\right)$}
%
\newcommand{\ds}{\displaystyle}
%
\renewcommand{\d}{\mathrm{\,d}}
\renewcommand{\i}{\mathrm{\,i\,}}
\newcommand{\e}{\mathrm{\,e\,}}
%%%%%%%%%%%%%%%%%%%%%%%%%%%%%%%%%%%%%%%%%%%


\begin{document}

%---------------------------------------------------------
% Réglage entête, pied de page
%---------------------------------------------------------
\rhead{\textbf{A. P{}. M. E. P{}.}}
\lhead{\small Corrigé}
\lfoot{\small{Polynésie}}
\rfoot{\small{26 juin 2026}}
\pagestyle{fancy}

\thispagestyle{empty}


%---------------------------------------------------------
% Titre
%---------------------------------------------------------
\begin{center}
	{\Large \bfseries
		\decofourleft~Diplôme National du Brevet~\decofourright\\[10pt]
		Série générale -- Polynésie -- 26 juin 2026
		\marginpar{\rotatebox{90}{\textbf{A. P{}. M. E. P{}.}}}\\[10pt]
		Corrigé }
\end{center}

\section*{\centering PREMIÈRE PARTIE : AUTOMATISMES (6 pts)}

\subsection*{Question 1}
	On range les valeurs de la série dans l'ordre croissant :

	\centerline{\begin{tblr}{colspec={*{7}{X[c,1cm]}},hline{1,2}={4}{solid},vline{4,5}={solid}}
		7&7&9&9&12&23&25\\
	\end{tblr}}

	Comme l'effectif total est impair (il y a 7 valeurs), la médiane est la valeur centrale (la quatrième), dans l'ordre croissant.

	La médiane de la série est 9.

\subsection*{Question 2}
	Le 4 qui est le premier chiffre non nul est le quatrième chiffre après la virgule, le chiffre des dix-millièmes.

	On a: \quad $\np{0,000457} = 4,57 \times \np{0,0001} = 4,57 \times 10^{-4}$.

	La notation scientifique de $\np{0,000457}$ est : \quad $4,57 \times 10^{-4}$.

\subsection*{Question 3}
	Comme le triangle est rectangle en B, on choisit [AB] comme base du triangle et donc, la hauteur correspondante est [BC].

	On a donc : \quad $\mathcal{A}_{\mathrm{ABC}} = \dfrac{\text{base} \times \text{hauteur}}{2} = \dfrac{\mathrm{AB} \times \mathrm{BC}}{2} = \dfrac{8 \times 6}{2} = \dfrac{48}{2} = 24$.

	Comme les longueurs sont données en centimètres, le calcul précédent donne une aire en centimètres carrés.

	L'aire du triangle ABC est donc \quad \np[cm^2]{24}.

\subsection*{Question 4}
	Puisque l'on pioche un beignet au hasard et qu'ils semblent identiques, on est en situation d'équiprobabilité.

	La probabilité d'un évènement est le nombre d'issues favorables à l'évènement, divisé par le nombre total d'issues.

	Il y a en tout :\quad $6 + 5 + 4 = 15$ beignets.

	Donc la probabilité de piocher un beignet à la framboise est :\quad $\dfrac{4}{15}$.

\subsection*{Question 5}
	Le coefficient multiplicateur d'une baisse de 10\,\% est : \quad $c = 1 - \dfrac{10}{100} = 1 - 0,1 = 0,9$.

 $800 \times 0,9 = 80 \times 9 =  720$, donc le prix, après réduction, est de: 720\,\euro{}.

	\emph{Variante :} $800 \times \dfrac{10}{100} = 80$, \quad 10\,\% de 800\,\euro{} représentent 80\,\euro{}.

$800 - 80 = 720$, donc le prix, après réduction, est de: 720\,\euro{}.

\subsection*{Question 6}
	Développons : \quad $B = 4y(3y - 1) = 4y \times 3y - 4y \times 1 = 4 \times 3 \times y \times y - 4y = 12y^2 - 4y$.

	La forme développée de $B$ est :\quad $B = 12y^2 - 4y$.

\subsection*{Question 7}
	On sait qu'un litre correspond à \np[dm^3]{1}, et que $\np[dm]{1} = \np[cm]{10}$, donc $\np[dm^3]{1} = 10^3\,\mathrm{cm}^3 = \np[cm^3]{1000}$.


	On peut faire un tableau de conversion :\quad
	\begin{tblr}{colspec={*{6}{X[c,4mm]}},colsep=1pt,hlines, vlines,vline{4}={1.2pt}}
		\SetCell[c=3]{c}{$\mathrm{dm}^3$} &  & & \SetCell[c=3]{c}$\mathrm{cm}^3$ \\
		&&3&5&7&0\\
	\end{tblr}

	\np[L]{3,57} correspond donc à \np[cm^3]{3570}.

	\textbf{Bonne réponse : D}.

\subsection*{Question 8}
	L'image de 4 par la fonction $f$ est :\quad $f(4) = 3 \times 4 - 5 = 12 - 5 = 7$.

	On a $f(4) = 7$.

	\textbf{Bonne réponse : B}.

\subsection*{Question 9}
	ABCD est un quadrilatère.

	Si on appelle O le point d'intersection des diagonales, les codages de la figure indiquent que DO = OB, et AO = OC.

	Ainsi, les diagonales se coupent en leurs milieux : cela garantit que ABCD est un parallélogramme.

	Comme le codage pour AO n'est pas le même que celui pour DO, on ne peut pas affirmer AO = DO : donc les diagonales n'ont pas forcément la même longueur, on ne peut pas affirmer que ABCD est un rectangle.

	Comme il n'y a aucun codage d'angle droit entre les diagonales, on ne peut pas affirmer que ABCD est un losange.

	Comme un carré est à la fois un losange et un rectangle, on ne peut pas affirmer que ABCD est un carré.

	\emph{Variante :} Pour l'absence de codage pour un rectangle, on peut aussi dire qu'il n'y a pas de codage d'angle droit entre deux côtés consécutifs, et pour l'absence de codage pour un losange, on peut aussi dire qu'il n'y a pas de codage de côtés consécutifs de même longueur.

	\medskip

	Quoi qu'il en soit, la seule nature qui est confirmée par les codages, c'est que le quadrilatère ABCD est un parallélogramme.

	\textbf{Bonne réponse : D}.

\needspace{10\baselineskip}
\section*{\centering DEUXIÈME PARTIE (14 pts)}

\section*{Exercice 1 (4 points)}
\subsection*{Partie A}
\begin{enumerate}
	\item \begin{enumerate}
		\item D'après la figure, on a : \quad $\mathrm{AB} = 3x + 1$.

		Avec $x = 10$, cela donne : \quad $\mathrm{AB} = 3 \times 10 + 1 = 31$.

		La longueur AB vaut donc 31.

		\item De même, on a :\quad $\mathrm{AD} = x - 2 = 10 - 2 = 8$.

		Comme ABCD est un rectangle, ses côtés opposés sont de même longueur, donc :

		$\mathcal{P}_{\mathrm{ABCD}} = 2\big(\mathrm{AB} + \mathrm{AD}\big) = 2\times (31 + 8) = 2\times 39 = 78$.

		Le périmètre de ABCD est bien de 78, pour $x=10$.
	\end{enumerate}
	\item En généralisant :\quad
	$\mathcal{P}_{\mathrm{ABCD}} = 2\big(\mathrm{AB} + \mathrm{AD}\big) = 2\times (3x + 1 + x - 2) = 2\times (4x - 1) = 8x - 2$.

	Le périmètre de ABCD est bien donné par l'expression \quad $8x-2$.
\end{enumerate}

\subsection*{Partie B}
\begin{enumerate}
	\item Si on choisit $x = 7$:
	\begin{itemize}
		\item le périmètre de ABCD est :\quad $8 \times 7 - 2 = 56 - 2 = 54$;
		\item le périmètre de IJK est : \quad $6 \times 7 + 9 = 42 + 9 = 51$.
	\end{itemize}
	Comme les deux périmètres ne sont pas égaux, le programme va afficher \og Les deux périmètres ne sont pas égaux \fg{} pendant 2 secondes.

	\item \begin{enumerate}
		\item Dans la cellule \textsf{B2}, on veut afficher le périmètre de ABCD, donné par l'expression $8x - 2$, en utilisant la valeur $x$ qui est renseignée dans la cellule \textsf{A2}.

		La formule correcte est donc :\quad \textsf{=8*A2-2}.

		Pourquoi les autres propositions sont fausses :
		\begin{itemize}
			\item \textsf{=8*5-2} \quad renvoie la valeur correcte pour la cellule \textsf{B2}, mais renverra toujours 38 quand on l'étirera vers le bas, car le \textsf{5} ne sera pas modifié;
			\item \textsf{=8*B2-2} \quad provoquera une erreur : c'est ce que l'on appelle une référence circulaire. La valeur de la cellule \textsf{B2} se fait à partir de la cellule \textsf{B2} : ça tourne en rond.
			\item \textsf{=8*A1-2} \quad provoquera une erreur : dans la cellule \textsf{A1}, il y a la lettre $x$, mais du coup, ce n'est pas un nombre que le tableur peut multiplier par 8.
		\end{itemize}

		\item Nour a constaté que pour $x = 5$ le périmètre de ABCD est \textbf{inférieur} à celui de IJK, par contre, pour $x = 6$, le périmètre de ABCD est \textbf{supérieur} à celui de IJK.

		Elle se dit donc que les courbes représentant les périmètres en fonction de $x$ doivent se croiser, entre l'abscisse 5 et l'abscisse 6.
	\end{enumerate}
	\item Résolvons l'équation :
	\begin{align*}
		8x - 2  & = 6x + 9        \\
		8x - 6x & = 9 + 2         \\
		2x      & = 11            \\
		x       & = \dfrac{11}{2}
	\end{align*}

	L'équation a une seule solution :\quad $\dfrac{11}{2} = 5,5$.

	C'est pour $x = 5,5$ que les deux figures ont le même périmètre.
\end{enumerate}

\needspace{10\baselineskip}
\section*{Exercice 2 (4 points)}
\begin{enumerate}
	\item \begin{enumerate}
		\item D'après le codage de la figure, le triangle CBF est isocèle en B, et donc les angles à la base ont la même mesure.

		On a donc, en degré:\quad $\widehat{\mathrm{CBF}} = \widehat{\mathrm{BCF}} = 74$.

		\item La somme des angles d'un triangle est un angle plat, de mesure 180\degres.

		On a donc, en degré:\quad $\widehat{\mathrm{CBF}}+\widehat{\mathrm{BCF}}+\widehat{\mathrm{CFB}} = 180$.

		En remplaçant par les mesure connues :\quad $\widehat{\mathrm{CBF}} + 74 + 74 = 180$.

		Finalement, on a donc : \quad $\widehat{\mathrm{CBF}} = 180 - (74 + 74) = 180 - 148 = 32$.

		L'angle $\widehat{\mathrm{CBF}}$ a donc une mesure de 32\degres.
	\end{enumerate}

	\item Dans le triangle ABC, on connaît les longueurs des trois côtés :
	\begin{itemize}
		\item le carré du côté le plus long est :\quad $\mathrm{AC}^2 = 5^2 = 25$;
		\item la somme des carrés des deux petits côtés est :\quad $\mathrm{AB}^2 + \mathrm{BC}^2 = 3^2 + 4^2 = 25$.
	\end{itemize}

	On constate que \quad $\mathrm{AC}^2 = \mathrm{AB}^2 + \mathrm{BC}^2$,\quad d'après la réciproque du théorème de Pythagore, on peut affirmer que le triangle ABC est un triangle rectangle en B.

	\item Le triangle BGF est rectangle en G, donc on peut utiliser les formules de trigonométrie, dans ce triangle.

	On s'intéresse à l'angle $\widehat{\mathrm{FBG}}$.

	\begin{itemize}
		\item On connaît la longueur en cm de l'hypoténuse du triangle :\quad $\mathrm{BF} = 4$;
		\item On connaît la longueur en cm du côté adjacent à l'angle $\widehat{\mathrm{FBG}}$:\quad $\mathrm{BG} = 2$;
	\end{itemize}

	On peut donc calculer le cosinus de l'angle :\quad $\cos\left(\widehat{\mathrm{FBG}}\right) = \dfrac{\text{côté adjacent}}{\text{hypoténuse}} = \dfrac{\mathrm{BG}}{\mathrm{BF}} = \dfrac{2}{4} = \dfrac{1}{2}$.

	On a donc : \quad $\widehat{\mathrm{FBG}} = \arccos\left(\dfrac{1}{2}\right) = 60\degres$.

	\item Les points A, B et G sont alignés si et seulement si l'angle $\widehat{\mathrm{ABG}}$ est un angle plat, dont la mesure est égale à 180\degres.

	Ici on a : \quad $\widehat{\mathrm{ABG}} = \widehat{\mathrm{ABC}} + \widehat{\mathrm{CBF}}+\widehat{\mathrm{FBG}}$.

	Avec les mesures connues en degré : \quad $\widehat{\mathrm{ABG}} = 90 + 32 + 60 = 182$.

	L'angle $\widehat{\mathrm{ABG}}$ n'est pas un angle plat, les points ne sont donc pas alignés.
\end{enumerate}

\needspace{10\baselineskip}
% saut de page si moins de 10 lignes disponibles
\section*{Exercice 3 (4 points)}
\begin{enumerate}
	\item On constate que Tahiti est à peu près au milieu entre l'équateur et le parallèle 30\degres{}S, et un peu plus à l'est que le méridien 150\degres{}O.

	Le décalage entre Tahiti et le méridien 150\degres{}O est comparable au décalage entre Los Angeles et le méridien 120\degres{}O, donc on peut dire que les coordonnées géographiques approximatives de Tahiti sont $(148\degres{}\mathrm{O} ; 15\degres{}\mathrm{S})$.

	\emph{Remarque :} L'île de Tahiti est en réalité aux coordonnées $(149\degres{}\mathrm{O} ; 17\degres{}\mathrm{S})$, mais la lecture sur le graphique ne permet pas d'avoir une précision au degré près, il y a donc une gamme de valeurs acceptables, comme réponse.

	\item Le trajet a duré 22 h 10, temps d'attente compris.

	Si on enlève 2 h 20 de temps d'attente, cela fait :

	$22\text{ h }10 - 2\text{ h }20 = 22\text{ h }10 - 2 \text{ h }10 - 0\text{ h }10 = 20\text{ h} - 0\text{ h }10 = 19\text{ h }50$.

	La durée combinée des deux vols est donc de 19 h 50 min.

	\item \begin{enumerate}
		\item Le diamètre de la médaille étant de \np[cm]{8,5}, son rayon est donc la moitié de cette longueur : \np[cm]{4,25}.

		En appliquant la formule rappelée au document 3 :

		$\mathcal{V}_{\text{médaille}} = \pi \times 4,25^2 \times 0,92 = \np{16,6175}\pi \approx52,205$.

		Comme toutes les longueurs étaient en centimètres, on a bien, au dixième près un volume de \np[cm^3]{52,2}.

		\item La masse d'argent sera donc proportionnelle à son volume. En faisant le calcul avec le volume arrondi au dixième obtenu précédemment :\quad
		$10,5 \times 52,2 = 548,1$.

		D'après ce calcul, la médaille contient donc environ \np[g]{548} d'argent, au gramme près.

		\emph{Remarque :} En réalité, les médailles d'argent des jeux de Paris 2024 sont faites d'argent (pureté 925/1000), pour environ 507 g, et de 18 g de fer issu des rénovations de la Tour Eiffel, pour un total de 525 g.

		Notre calcul donne un résultat supérieur, car il ne prend pas en compte les motifs en creux de la médaille, qui font qu'il y a un peu moins d'argent que le cylindre plein.


	\end{enumerate}
\end{enumerate}
\end{document}